\paragraph{window-\(6\) hidden fiber 的 cubical 规范形、Bernoulli 因子化与 UV alias 超平面}
在定理 \ref{thm:xi-time-part9x-visible-condexp-finite-index-spectrum} 与定理 \ref{thm:xi-time-part9x-equivariant-completion-bit-lower-bound} 已经把 window-\(6\) 的 visible/hidden 裂分、局部可逆容量与对称 completion 溢价固定之后，还可仅凭 \(m=6\) 的显式审计数据，把 hidden fiber 的局部组合形、Fourier 零模与连续极限下的 alias 机制压缩为另一条有限尺度闭链。以下沿用
\[
\Omega_6:=\{0,\dots,63\},
\qquad
\pi:=\Fold_6^{\mathrm{bin}}:\Omega_6\to X_6,
\qquad
F_w:=\pi^{-1}(w),
\qquad
d(w):=|F_w|,
\]
并采用全文的稳定词记号
\[
X_m:=\{u\in\{0,1\}^m:\ u_i u_{i+1}=0\ \ (1\le i<m)\}.
\]

\begin{theorem}[window-\texorpdfstring{$6$}{6} 最小退化扇区的 rigid \texorpdfstring{$X_4$}{X4} 子覆盖]\label{thm:xi-time-part9xab-window6-minsector-rigid-x4-subcover}
\leanverified{paper\_xi\_time\_part9xab\_window6\_minsector\_rigid\_x4\_subcover}
定义
\[
S_{6,2}:=\{w\in X_6:\ d(w)=2\},
\qquad
\iota_4:X_4\hookrightarrow X_6,
\qquad
\iota_4(v):=(v,0,1).
\]
则有精确恒等式
\[
S_{6,2}=\iota_4(X_4).
\]
换言之，window-\(6\) 的全部最小退化态恰由后缀固定为 \(01\) 的稳定词组成，并与 \(X_4\) 逐词同构。
\end{theorem}
\begin{proof}
由定理 \ref{thm:xi-foldbin-zeckendorf-monotone-staircase-law} 在 \(m=6\) 时的三档公式，
\[
d(w)=2
\quad\Longleftrightarrow\quad
w_6=1.
\]
若 \(w\in X_6\) 满足 \(w_6=1\)，则稳定性约束 \(w_5w_6=0\) 强制 \(w_5=0\)，故 \(w\) 唯一写成 \(w=(v,0,1)\)，其中前四位 \(v\in X_4\)。反之，任取 \(v\in X_4\)，由于 \(v_4\cdot 0=0\) 且 \(0\cdot 1=0\)，故 \((v,0,1)\) 仍属 \(X_6\)。于是
\[
\{w\in X_6:\ w_6=1\}=\iota_4(X_4).
\]
与上式合并即得
\[
S_{6,2}=\iota_4(X_4).
\]
证毕。
\end{proof}

\begin{theorem}[window-\texorpdfstring{$6$}{6} hidden fiber 的仿射 Boolean 方块标准形]\label{thm:xi-time-part9xab-window6-hidden-fiber-affine-boolean-normal-form}
\leanverified{paper\_xi\_time\_part9xab\_window6\_hidden\_fiber\_affine\_boolean\_normal\_form}
定义
\[
\beta:\{0,1\}^2\to \ZZ,
\qquad
\beta(a,b):=21a+34b,
\]
并令
\[
I_2:=\{(0,0),(0,1)\},
\qquad
I_3:=\{(0,0),(1,0),(0,1)\},
\qquad
I_4:=\{0,1\}^2.
\]
若对任意 \(J\subseteq \{0,1\}^2\) 记
\[
\beta(J):=\{\beta(u):\ u\in J\},
\]
则对每个 \(w\in X_6\) 都有
\[
F_w=V_6(w)+\beta(I_{d(w)}).
\]
特别地，
\[
\beta(I_2)=\{0,34\},
\qquad
\beta(I_3)=\{0,21,34\},
\qquad
\beta(I_4)=\{0,21,34,55\},
\]
并且 \(I_2,I_3,I_4\) 都是 Boolean 方块 \(\{0,1\}^2\) 在乘积序下的序理想。
\end{theorem}
\begin{proof}
由定理 \ref{thm:xi-foldbin-zeckendorf-monotone-staircase-law} 的三档分类与表 \ref{tab:fold6_bin_uplift_choice_collapse}，
\[
F_w=
\begin{cases}
\{V_6(w),\,V_6(w)+34\}, & d(w)=2,\\[4pt]
\{V_6(w),\,V_6(w)+21,\,V_6(w)+34\}, & d(w)=3,\\[4pt]
\{V_6(w),\,V_6(w)+21,\,V_6(w)+34,\,V_6(w)+55\}, & d(w)=4.
\end{cases}
\]
又因为
\[
55=21+34=\beta(1,1),
\]
故三类偏移集合恰分别等于 \(\beta(I_2),\beta(I_3),\beta(I_4)\)，从而
\[
F_w=V_6(w)+\beta(I_{d(w)}).
\]
最后，在乘积序 \((a,b)\le (a',b')\iff a\le a',\ b\le b'\) 下，
\[
I_2=\{u\in\{0,1\}^2:\ u\le (0,1)\},
\]
\[
I_3=\{u\in\{0,1\}^2:\ u\le (1,0)\ \text{或}\ u\le (0,1)\},
\qquad
I_4=\{u\in\{0,1\}^2:\ u\le (1,1)\},
\]
均为下闭集，故皆为序理想。证毕。
\end{proof}

\begin{corollary}[任意有限体积下 hidden fiber 的全局可缩性]\label{cor:xi-time-part9xab-window6-hidden-fiber-global-contractibility}
对任意有限指标集 \(\Lambda\)，定义
\[
\pi_\Lambda:=\pi^\Lambda:\Omega_6^\Lambda\to X_6^\Lambda.
\]
固定 \(w=(w_x)_{x\in\Lambda}\in X_6^\Lambda\)，并定义 cubical realization
\[
\mathfrak F(w):=\prod_{x\in\Lambda} I_{d(w_x)}
\subset ([0,1]^2)^\Lambda.
\]
则 \(\pi_\Lambda^{-1}(w)\) 与 \(\mathfrak F(w)\) 的顶点集存在规范双射，而且 \(\mathfrak F(w)\) 可缩。因而
\[
\widetilde H_k(\mathfrak F(w);\ZZ)=0
\qquad (k\ge 0).
\]
\end{corollary}
\begin{proof}
由定理 \ref{thm:xi-time-part9xab-window6-hidden-fiber-affine-boolean-normal-form}，
\[
\pi_\Lambda^{-1}(w)
=
\prod_{x\in\Lambda}\bigl(V_6(w_x)+\beta(I_{d(w_x)})\bigr).
\]
由于 \(\beta\) 在 \(I_2,I_3,I_4\) 上均为单射，故每个 \(n_x\in V_6(w_x)+\beta(I_{d(w_x)})\) 唯一对应于某个顶点 \(u_x\in I_{d(w_x)}\)。逐坐标合并，得到 \(\pi_\Lambda^{-1}(w)\) 与 \(\mathfrak F(w)\) 顶点集之间的规范双射。

再者，把 \(I_2,I_3,I_4\) 视为 \([0,1]^2\) 中的有限 cubical complex，则 \(I_2\) 是闭区间，\(I_3\) 是一个 \(L\) 形树，\(I_4\) 是单位正方形，三者都可缩。有限乘积仍可缩，因此 \(\mathfrak F(w)\) 可缩。可缩空间的约化同调全消失，故
\[
\widetilde H_k(\mathfrak F(w);\ZZ)=0
\qquad (k\ge 0).
\]
证毕。
\end{proof}

\begin{theorem}[四点纤维的 rank-two Bernoulli 因子化]\label{thm:xi-time-part9xab-window6-fourpoint-fiber-ranktwo-bernoulli-factorization}
\leanverified{paper\_xi\_time\_part9xab\_window6\_fourpoint\_fiber\_ranktwo\_bernoulli\_factorization}
对任意 \(w\in X_6\)，定义局部偏移 Laplace 变换
\[
L_w(t):=\frac1{d(w)}
\sum_{n\in F_w} e^{\,t(n-V_6(w))}.
\]
则有
\[
L_w(t)=\frac{1+e^{34t}}{2}
\qquad (d(w)=2),
\]
\[
L_w(t)=\frac{1+e^{21t}+e^{34t}}{3}
\qquad (d(w)=3),
\]
\[
L_w(t)=\frac{1+e^{21t}+e^{34t}+e^{55t}}{4}
=
\frac{1+e^{21t}}2\cdot \frac{1+e^{34t}}2
\qquad (d(w)=4).
\]
因此当 \(d(w)=4\) 时，均匀 hidden 偏移分布与
\[
21B_1+34B_2,
\qquad
B_1,B_2\stackrel{\mathrm{i.i.d.}}{\sim}\mathrm{Bernoulli}\!\left(\tfrac12\right),
\]
完全同分布。
\end{theorem}
\begin{proof}
由定理 \ref{thm:xi-time-part9xab-window6-hidden-fiber-affine-boolean-normal-form}，
\[
F_w-V_6(w)=
\begin{cases}
\{0,34\}, & d(w)=2,\\
\{0,21,34\}, & d(w)=3,\\
\{0,21,34,55\}, & d(w)=4.
\end{cases}
\]
把这三种偏移集合代入 \(L_w(t)\) 的定义，即得前三个闭式。

当 \(d(w)=4\) 时，
\[
\{0,21,34,55\}=\{0,21\}+\{0,34\},
\]
并且四个和都恰好出现一次。因此均匀分布于 \(\{0,21,34,55\}\) 的随机变量，与两个独立公平 Bernoulli 变量经线性映射
\[
(a,b)\longmapsto 21a+34b
\]
的像分布完全一致。Laplace 变换遂分解为
\[
\frac{1+e^{21t}}2\cdot \frac{1+e^{34t}}2.
\]
证毕。
\end{proof}

\begin{theorem}[window-\texorpdfstring{$6$}{6} 局部 Fourier 乘子的精确零模集]\label{thm:xi-time-part9xab-window6-local-fourier-multipliers-exact-zero-set}
\leanverified{paper\_xi\_time\_part9xab\_window6\_local\_fourier\_multipliers\_exact\_zero\_set}
令
\[
\omega:=e^{2\pi i/64}.
\]
对三类局部模板定义归一化 Fourier 乘子
\[
m_2(t):=\frac{1+\omega^{34t}}2,
\qquad
m_3(t):=\frac{1+\omega^{21t}+\omega^{34t}}3,
\]
\[
m_4(t):=\frac{1+\omega^{21t}+\omega^{34t}+\omega^{55t}}4
=
\frac{(1+\omega^{21t})(1+\omega^{34t})}{4},
\qquad
t\in \ZZ/64\ZZ.
\]
则其零点集精确为
\[
Z(m_2)=\{16,48\},
\qquad
Z(m_3)=\varnothing,
\qquad
Z(m_4)=\{16,32,48\}.
\]
\end{theorem}
\begin{proof}
先看 \(m_2\)。有
\[
m_2(t)=0
\quad\Longleftrightarrow\quad
\omega^{34t}=-1
\quad\Longleftrightarrow\quad
34t\equiv 32 \pmod{64}.
\]
两边同除以 \(2\)，得到
\[
17t\equiv 16 \pmod{32}.
\]
由于 \(17^{-1}\equiv 17\pmod{32}\)，故
\[
t\equiv 16 \pmod{32},
\]
于是
\[
Z(m_2)=\{16,48\}.
\]

再看 \(m_4\)。由因子化公式，
\[
m_4(t)=0
\quad\Longleftrightarrow\quad
\omega^{21t}=-1\ \text{或}\ \omega^{34t}=-1.
\]
第二种情形已给出 \(t\in\{16,48\}\)。第一种情形等价于
\[
21t\equiv 32 \pmod{64}.
\]
由于 \(21^{-1}\equiv 61\pmod{64}\)，遂得
\[
t\equiv 61\cdot 32\equiv 32 \pmod{64}.
\]
故
\[
Z(m_4)=\{16,32,48\}.
\]

最后考察 \(m_3\)。若 \(m_3(t)=0\)，则存在单位复数
\[
\zeta:=\omega^{21t},
\qquad
\eta:=\omega^{34t}
\]
满足
\[
1+\zeta+\eta=0.
\]
于是 \(\eta=-(1+\zeta)\)，从而
\[
1=|\eta|^2=|1+\zeta|^2=2+2\Re(\zeta).
\]
故
\[
\Re(\zeta)=-\frac12.
\]
这迫使 \(\zeta\) 为本原三次单位根。然而 \(\zeta\) 同时又是 \(64\) 次单位根，其阶数必须整除 \(64\)，与 \(3\nmid 64\) 矛盾。因此 \(m_3\) 无零点，即
\[
Z(m_3)=\varnothing.
\]
证毕。
\end{proof}

\begin{theorem}[全局 hidden averaging 的 alias 超平面]\label{thm:xi-time-part9xab-window6-hidden-averaging-alias-hyperplanes}
固定有限指标集 \(\Lambda\) 与可见构型 \(w=(w_x)_{x\in\Lambda}\in X_6^\Lambda\)。把每个全局原像 \((n_x)_{x\in\Lambda}\in \pi_\Lambda^{-1}(w)\) 唯一写成
\[
n_x=V_6(w_x)+\delta_x,
\qquad
\delta_x\in \beta(I_{d(w_x)}).
\]
定义 offsets 上的均匀平均算子
\[
E_w f
:=
\frac{1}{\prod_{x\in\Lambda} d(w_x)}
\sum_{(\delta_x)_x\in \prod_{x\in\Lambda}\beta(I_{d(w_x)})}
f\bigl((\delta_x)_{x\in\Lambda}\bigr).
\]
对任意 Fourier 角色
\[
\chi_t(\delta):=
\exp\!\Bigl(\frac{2\pi i}{64}\sum_{x\in\Lambda} t_x\delta_x\Bigr),
\qquad
t=(t_x)_{x\in\Lambda}\in (\ZZ/64\ZZ)^\Lambda,
\]
有
\[
E_w\chi_t=\prod_{x\in\Lambda} m_{d(w_x)}(t_x),
\]
其中 \(m_2,m_3,m_4\) 由定理 \ref{thm:xi-time-part9xab-window6-local-fourier-multipliers-exact-zero-set} 定义。因而精确零模集合
\[
\mathcal Z_w:=\{t\in (\ZZ/64\ZZ)^\Lambda:\ E_w\chi_t=0\}
\]
满足
\[
\mathcal Z_w
=
\bigcup_{x:\,d(w_x)=2}\{t:\ t_x\in\{16,48\}\}
\ \cup\
\bigcup_{x:\,d(w_x)=4}\{t:\ t_x\in\{16,32,48\}\}.
\]
特别地，三点纤维不贡献任何精确零超平面。
\end{theorem}
\begin{proof}
由定理 \ref{thm:xi-time-part9xab-window6-hidden-fiber-affine-boolean-normal-form}，全局 hidden fiber 按坐标分解为直积
\[
\pi_\Lambda^{-1}(w)-V_6(w)
=
\prod_{x\in\Lambda}\beta(I_{d(w_x)}),
\]
其中 \(V_6(w):=(V_6(w_x))_{x\in\Lambda}\)。因此对角色 \(\chi_t\) 的均匀平均可完全分离：
\[
\begin{aligned}
E_w\chi_t
&=
\frac{1}{\prod_{x\in\Lambda} d(w_x)}
\sum_{(\delta_x)_x\in \prod_{x\in\Lambda}\beta(I_{d(w_x)})}
\prod_{x\in\Lambda}
\exp\!\Bigl(\frac{2\pi i}{64}t_x\delta_x\Bigr)\\
&=
\prod_{x\in\Lambda}
\left(
\frac1{d(w_x)}
\sum_{\delta_x\in \beta(I_{d(w_x)})}
e^{2\pi i t_x\delta_x/64}
\right)\\
&=
\prod_{x\in\Lambda} m_{d(w_x)}(t_x).
\end{aligned}
\]
于是 \(E_w\chi_t=0\) 当且仅当某个坐标因子为零。再将定理 \ref{thm:xi-time-part9xab-window6-local-fourier-multipliers-exact-zero-set} 的精确零点集逐坐标代入，即得
\[
\mathcal Z_w
=
\bigcup_{x:\,d(w_x)=2}\{t:\ t_x\in\{16,48\}\}
\ \cup\
\bigcup_{x:\,d(w_x)=4}\{t:\ t_x\in\{16,32,48\}\}.
\]
而 \(Z(m_3)=\varnothing\)，故三点纤维不贡献任何精确零超平面。证毕。
\end{proof}

\begin{conjecture}[sitewise window-\texorpdfstring{$6$}{6} lift 的 UV alias 原理]\label{conj:xi-time-part9xab-window6-lift-uv-alias-principle}
沿用猜想 \ref{conj:xi-time-part9x-finite-index-symmetry-premium-visible-gff} 的设定，设某一临界六顶点系综中 visible 高度场在适当归一化下收敛到标量高斯自由场。则任一 canonical sitewise window-\(6\) lift 的 hidden averaging 在连续极限中只可能留下两类效应：
\[
\text{(i) 有限局域熵势或接触项重整化},\qquad
\text{(ii) 仅支撑于 } \{16,32,48\}\text{ 的 UV alias 缺口}.
\]
更精确地，hidden averaging 不会创造新的红外零模，也不会改变 visible sector 所决定的极限 Green 核；它在离散频域中的全部精确零模，均由定理 \ref{thm:xi-time-part9xab-window6-hidden-averaging-alias-hyperplanes} 的坐标超平面生成，而三点纤维对应的局部模板在频域上保持严格零缺口。
\end{conjecture}
\begin{proof}[证据]
推论 \ref{cor:xi-time-part9xab-window6-hidden-fiber-global-contractibility} 表明，任意有限体积上的 hidden fiber 都可由若干可缩 cubical 因子的乘积实现，从而纤维本体不携带新的局域高阶拓扑。定理 \ref{thm:xi-time-part9xab-window6-fourpoint-fiber-ranktwo-bernoulli-factorization} 又把全部四点纤维精确裂解为两个独立 Bernoulli 模，而定理 \ref{thm:xi-time-part9xab-window6-local-fourier-multipliers-exact-zero-set} 说明：局部精确零模只可能出现在 \(d=2\) 与 \(d=4\) 模板上，频率严格限制在
\[
\{16,48\}\quad\text{与}\quad \{16,32,48\},
\]
其中 \(d=3\) 模板完全零自由。定理 \ref{thm:xi-time-part9xab-window6-hidden-averaging-alias-hyperplanes} 进一步把这一现象推广到任意有限体积：全部精确零模只是坐标方向上的 alias 超平面，与体积增长无关，也不产生新的低频连续族。

因此，在猜想 \ref{conj:xi-time-part9x-finite-index-symmetry-premium-visible-gff} 所要求的 GFF 归一化下，hidden averaging 的频域效应应当只能停留在离散 UV 层，并在连续极限中塌缩为局域接触型修正；真正控制红外质量零模的仍是 visible sector。证毕。
\end{proof}

\noindent
由此，window-\(6\) 的 hidden sector 不仅在纤维计数上是有限的，而且在局部组合与频域上都被完全 rigidify：最小二重点层恰是 \(X_4\) 的 rigid 后缀子覆盖；任一局部纤维都由二维 Boolean 方块的仿射序理想给出；四点纤维裂解为两条独立 Bernoulli 模，而全部精确频域零模又只落在 \(\{16,32,48\}\) 的 alias 超平面上。于是，若 visible sector 的高斯自由场极限成立，则 hidden averaging 的剩余作用便应当只能停留在局域 UV 修正层面。

\endinput
